% why-this-physics.tex — Section 5: Why This Physics?

\section{Why This Physics?}
\label{sec:why-this-physics}

\subsection{General relativity from locality}
\label{subsec:gr}

Paper~6~\citep{ehrlich2026gr} argues that self-modeling locality on a
lattice provides a UV completion for entanglement-based gravity.  The
argument has two layers: a UV layer (derived from self-modeling) and a
bridge to Einstein's equations (inherited from Jacobson's
thermodynamic route~\citep{jacobson1995}).

\textbf{UV layer.}  A self-modeler embedded in a lattice interacts
with neighbors via a SWAP-type Hamiltonian (the minimal Hamiltonian
that preserves the body-model factorization).  The resulting dynamics
satisfy a Lieb-Robinson bound~\citep{liebrobinson1972}, giving an
emergent causal structure.  Entanglement across any bipartition of the
lattice obeys an area law~\citep{bombelli1986, bekenstein1973}: the
entanglement entropy of a region scales with its boundary area, not
its volume.  This is forced by the locality of the Hamiltonian and the
spectral gap, not assumed.

\textbf{Jacobson bridge.}  Jacobson~\citep{jacobson1995} showed that
any theory satisfying (i) a local entropy-area relation $\delta S =
\eta\,\delta A$ for causal horizons, (ii) local Lorentz invariance,
and (iii) the Clausius relation $\delta Q = T\,\delta S$ for local
Rindler observers produces Einstein's field equations as an equation
of state.  The UV layer provides ingredient (i).  Ingredients (ii) and
(iii) are inherited as standard assumptions.

\textbf{Four standard gaps.}  The derivation is conditional on four
gaps, none of which are bespoke to this program:
\begin{enumerate}[nosep]
\item \emph{Emergent Lorentz invariance}: the lattice must produce
  continuous Lorentz symmetry in the continuum limit.
\item \emph{Lattice Bisognano-Wichmann}: the Unruh effect must hold
  for lattice Rindler observers, not just in the continuum.
\item \emph{Local equilibrium}: the Clausius relation must hold
  locally, requiring fast thermalization near causal horizons.
\item \emph{Continuum limit}: the lattice theory must have a
  well-defined continuum limit producing smooth spacetime.
\end{enumerate}
All four are shared by every lattice quantum gravity program.  They
are identified honestly; none are claimed to be closed.

\subsection{The Standard Model from $\hthree$}
\label{subsec:sm}

Paper~7~\citep{ehrlich2026sm} argues that the Standard Model gauge
group emerges from the exceptional Jordan algebra $\hthree$, the
unique non-composable self-modeling structure.  The argument is
structural, not spectral-theoretic.

\textbf{Why $\hthree$.}  Self-modeling forces Jordan algebras (Step~3
of the chain).  The question is which Jordan algebra governs the
universe as a whole.  The answer follows from a purely algebraic
argument whose key premise is radical relativity itself.

\begin{claim}[Non-composability]
\label{claim:non-comp}
The universe's self-modeling algebra is $\hthree$.
\end{claim}

\textbf{Argument.}
\begin{enumerate}[nosep]
\item Self-modeling forces a Jordan algebra
  (Paper~5~\citep{ehrlich2026qm}, theorem).
\item Special (composable) Jordan algebras admit composites with
  nontrivial systems: for any special $J$, a well-behaved composite
  $J \otimes K$ exists~\citep{bgw2020, hancheolsen1985}.
\item Under radical relativity (Postulate~\ref{post:l4}), all
  consistent mathematical structures exist.  If composites of $J$
  exist mathematically, they exist.  Therefore $J$ is a subsystem of
  those composites.
\item Define the \emph{universe} as the maximal self-modeling structure
  containing the observer.  (Maximality is well-defined because
  assumption A1 restricts to finite dimension, bounding the
  Jordan-algebraic dimension from above.)
\item If the universe's algebra $U$ were composable (special), then
  $U$ embeds in an associative $C^*$-algebra
  $A$~\citep{hancheolsen1985}.  Composites $A \otimes B$ exist for
  nontrivial $B$ (by steps 2--3), and the composite
  $A \otimes B \cong M_{nk}(\mathbb{C})$ is again a $C^*$-algebra.
  By Paper~5's main result~\citep{ehrlich2026qm}, any
  $\Mnsa{m}$ admits a faithful self-model (take $V_M = V_B$,
  $\varphi = \mathrm{id}$).  Therefore the composite carries its own
  faithful self-modeling structure, strictly larger than $U$, in which
  the observer is embedded.  This contradicts maximality of $U$.
\item Therefore the universe's algebra is non-composable
  (contrapositive of step~5).
\item The unique non-composable simple Jordan algebra is $\hthree$
  (JvNW classification~\citep{jvnw1934} plus
  Zel'manov~\citep{zelmanov1979}).
\end{enumerate}

\begin{remark}
Under radical relativity, algebraic composability \emph{is} the
physical question of subsystem-hood.  There is no separate ``physical''
layer in which composability might differ from the algebraic property.
This is why the argument uses no quantum entanglement or physical
reasoning --- only the Jordan-algebraic notion of composability, which
is prior to quantum mechanics in the derivation chain.  Step~7
requires simplicity (the universe's algebra is not a direct sum); this
is one of Paper~5's four standing assumptions (simple EJA via
trace-form non-degeneracy), an honest cost recorded in
Sec.~\ref{subsec:obj-gaps}.
\end{remark}

\textbf{The gauge group.}  $\mathrm{Aut}(\hthree) = \Ffour$.  An
observer, modeled as a $C^*$-subsystem
(Paper~5~\citep{ehrlich2026qm}), picks a rank-1 idempotent $E_{11}
\in \hthree$, inducing a Peirce decomposition into three sectors:
$V_1$ (observer, 1-dim), $V_{1/2}$ (interface, 16-dim), and $V_0$
(complement, 10-dim).  The stabilizer of $E_{11}$ in $\Ffour$ is
$\Spin{9}$.  A second choice --- a complex structure $u \in S^6
\subset \mathrm{Im}(\mathbb{O})$ --- breaks $\Ffour$ differently,
stabilizing $[\SU{3} \times \SU{3}] / \mathbb{Z}_3$.  The
intersection of these two stabilizers is the Standard Model gauge
group $\SMgauge$.  This is the Todorov-Drenska
theorem~\citep{todorovdrenska2018}.

\textbf{Chirality.}  The same complex structure $u$ that determines
the gauge group simultaneously determines the chiral representation.
The Clifford algebra $\mathrm{Cl}(6)$, generated by the six imaginary
octonions orthogonal to $u$, admits a Witt decomposition whose volume
form produces an automatic left-handed embedding of Standard Model
fermions.  One choice does three things: gauge group,
complexification, and left chirality.  This unifies the
Todorov-Drenska route (gauge group) with the Furey
route~\citep{furey2018} (chirality) via a single algebraic input.

\textbf{Gap C.}  The argument that the observer's $C^*$-algebra
nature induces complexification of the Peirce half-space ($\Spin{9}
\to \Spin{10}$) is physical reasoning, not a theorem.  This is Gap~C,
the one remaining real gap in the SM derivation.  Two independent
motivations exist (algebraic and thermodynamic), but neither
constitutes a proof.  The honest assessment: Gap~C may admit only an
evolutionary resolution --- ``no chirality $\to$ no electroweak
symmetry breaking $\to$ no mass hierarchy $\to$ no chemistry $\to$ no
self-modelers'' --- the same status as ``why eyes.''
Overwhelmingly plausible; not a theorem; perhaps not closeable as one
(Sec.~\ref{subsec:obj-gaps}).

\begin{claim}[Uniqueness of $\rhoJ$]
\label{claim:rhoJ}
Within $\hthree$, the heuristic density~\eqref{eq:rho} resolves
uniquely.  Paper~1~\citep{ehrlich2026measure} proves:
$\rhoJ(X) = \det(X) \cdot (\mathrm{Tr}(X^2) - \tfrac{1}{3})$ is the
unique lowest-degree $\Ffour$-invariant polynomial on the state space
of $\hthree$ that vanishes at rank-deficient states ($\det = 0$) and
at thermal equilibrium ($\mathrm{Tr}(X^2) = \tfrac{1}{3}$), and is
non-negative throughout.
\end{claim}

The cubic norm $\det(X) = \lambda_1 \lambda_2 \lambda_3$ is the
exceptional invariant of $\hthree$ --- it does not exist for ordinary
quantum mechanics ($h_n(\mathbb{C})$).  $\rhoJ > 0$ \emph{requires}
the exceptional structure.  The purity factor
$\mathrm{Tr}(X^2) - \tfrac{1}{3}$ measures distance from thermal
equilibrium.  Their product is the unique algebraic experiential
measure: (three-way correlation) $\times$ (thermodynamic
disequilibrium).

\begin{remark}[The cubic prediction]
\label{rem:cubic-prediction}
$\rhoJ$ depends on $\det(X)$, a cubic invariant that vanishes when
\emph{any} Peirce sector is empty.  Integrated information
$\Phi$~\citep{tononi2008} is quadratic (pairwise integration),
remaining nonzero when one sector is disrupted.  The prediction:
disrupting one of three sectors produces a sharp loss of self-modeling
quality ($\det \to 0$), not a gradual reduction.  This is the
program's most distinguishing falsifiable prediction.
\end{remark}

\subsection{Three generations from the Peirce tower [Speculative]}
\label{subsec:peirce-tower}

The Standard Model has three generations of fermions with a steep mass
hierarchy ($m_t / m_u \sim 10^5$).  Boyle~\citep{boyle2020} noted
that $\hthree$ contains three off-diagonal octonionic entries
$(x_1, x_2, x_3)$ permuted by $\mathrm{SO}(8)$ triality, providing a
structural reason for exactly three generations.  Triality treats all
three symmetrically and does not explain the mass hierarchy.

The iterated Peirce decomposition breaks this symmetry.  Decomposing
$\hthree$ under $E_{11}$ gives $V_1 + V_{1/2} + V_0$, where
$V_{1/2} = \mathbb{O}^2$ contains $x_2$ and $x_3$ but not $x_1$.
The complement $V_0 = h_2(\mathbb{O})$ is itself a Jordan algebra;
decomposing it under a rank-1 idempotent $E_{22}$ reveals $x_1$ as
the interface at the second level.  The tower bottoms out at the
third level ($h_1(\mathbb{O}) = \mathbb{R}$).  Interface dimensions
descend: 16, 8, 1.

The observer's existence breaks triality to a $2+1$ split: $x_2$ and
$x_3$ are direct interfaces (level~1), while $x_1$ is one level
removed inside the complement (level~2).  Coupling strength to the
observer --- and therefore mass --- decreases with Peirce depth.  The
third generation (heaviest) couples at level~1; the first generation
(lightest) couples at level~2.

\begin{remark}
This is explicitly speculative.  Whether the Peirce tower reproduces
the quantitative mass hierarchy (the actual ratio $m_t / m_u$)
depends on the spectral action on $\hthree$, which has not been
computed.  The structural prediction is qualitative: the hierarchy
exists and has a $2+1$ pattern.  See Sec.~\ref{subsec:open-problems}
for the status of the spectral action computation.
\end{remark}
